\documentclass[11pt]{scrartcl}
\usepackage[a4paper,margin=2.4cm]{geometry}
\usepackage{../common/rlstudy}

\title{\textbf{Week 22 — Extended reading}\\\large Constrained MDPs, CPO/FOCOPS/P3O, Lagrangian Safety}
\author{STUDY-RL}
\date{\today}

\begin{document}
\maketitle

\section*{Paper 1 --- ``Benchmarking Safe Exploration in Deep RL'' (PPO-Lagrangian)}
\textbf{Authors:} Alex Ray, Joshua Achiam, Dario Amodei. \textbf{Year:} 2019.
\textbf{Venue:} OpenAI tech.\ report. \textbf{arXiv:} 1910.01708.

\paragraph{Summary.} Introduces \emph{Safety Gym}, a suite of MuJoCo
robotics tasks with constraint signals (e.g.\ entering hazards, contacting
fragile objects), and benchmarks PPO-Lagrangian and CPO on them. The paper
formalises the CMDP API for deep RL, gives canonical hyperparameters, and
shows that simple primal-dual PPO already matches or beats CPO on most tasks
when carefully tuned --- a foundational empirical reference.

\paragraph{Why it matters for this week.} It is the canonical implementation
template for PPO-Lagrangian. Our \texttt{project/agent.py} follows the same
architecture (separate $V_r$, $V_c$ heads; combined advantage; projected
multiplier update) and its decision to keep the rollout small for the dual
loop comes directly from this paper.

\paragraph{Annotated key passage.}
\begin{quote}
``Despite the simplicity of the Lagrangian approach, getting it to work
reliably in practice required several engineering decisions: separate cost
critics, advantage standardisation \emph{only} on the reward, and a learning
rate for $\lambda$ that is one to two orders of magnitude smaller than the
policy learning rate to avoid oscillation.''
\end{quote}
\emph{Commentary.} Every implementer rediscovers this trio of tricks within
their first week of hacking on safe RL --- the paper saves you that week. The
small dual-LR is essentially a poor-person's PID damping; the same effect is
later formalised by Stooke et al.\ as PID-Lagrangian.

% -------------------------------------------------------------------
\section*{Paper 2 --- ``Constrained Policy Optimization'' (CPO)}
\textbf{Authors:} Joshua Achiam, David Held, Aviv Tamar, Pieter Abbeel.
\textbf{Year:} 2017. \textbf{Venue:} ICML. \textbf{arXiv:} 1705.10528.

\paragraph{Summary.} The first deep-RL algorithm with a per-iteration
near-feasibility guarantee. CPO extends TRPO by adding a cost-constraint to
the trust-region QP, deriving a closed-form dual solution (Theorem 4.1 in
the paper). It also describes a \emph{recovery step} when the linearised
constraint cannot be satisfied inside the trust region.

\paragraph{Why it matters.} It is the theoretical reference point for all
later constrained RL work; FOCOPS, P3O, and PID-Lag all benchmark against
it. Even if you choose not to implement Fisher-vector products on a single
GPU, you must understand CPO to read modern papers.

\paragraph{Annotated key passage.}
\begin{quote}
``Theorem 1 (CPO performance bound). For any policies $\pi'$ and $\pi$,
$J_c(\pi') \leq J_c(\pi) + \tfrac{1}{1-\gamma}\E_{s \sim d^\pi, a \sim \pi'}[A_c^\pi(s,a)]
 + \tfrac{\sqrt{2}\gamma\epsilon_c^\pi}{(1-\gamma)^2}\sqrt{\tfrac{1}{2}\bar D_\KL(\pi'\Vert\pi)}.$''
\end{quote}
\emph{Commentary.} This inequality is the central object of CPO. The
$1/(1-\gamma)^2$ scaling tells you the long-horizon penalty for trust-region
slack --- short-horizon tasks tolerate larger $\delta$. The
$\sqrt{\bar D_\KL}$ dependence is what makes the trust-region argument
non-vacuous, and exactly why TRPO's KL constraint pulls double duty here.

% -------------------------------------------------------------------
\section*{Paper 3 --- ``FOCOPS''}
\textbf{Authors:} Yiming Zhang, Quan Vuong, Keith Ross. \textbf{Year:} 2020.
\textbf{Venue:} NeurIPS. \textbf{arXiv:} 2002.06506.

\paragraph{Summary.} Solves the CPO problem in two phases: (i) compute the
non-parametric optimal policy in closed form (KL-regularised reward maximisation
with a cost multiplier); (ii) project onto the parametric family by KL
minimisation. Replaces all Fisher--vector products with first-order updates.

\paragraph{Why it matters.} Best-in-class fit to a single-GPU budget.
Conjugate-gradient steps inside CPO can dominate the wall-clock cost on
deep policies; FOCOPS removes them while retaining a CPO-style theoretical
worst-case bound (in expectation, not worst-case-per-state).

\paragraph{Annotated key passage.}
\begin{quote}
``Lemma 1. Given fixed multipliers $\nu, \lambda$, the optimal solution to
the relaxed problem in non-parametric policy space is $\pi^\star(a|s)
\propto \pi_k(a|s) \exp(\tfrac{1}{\lambda}(A_r - \nu A_c))$.''
\end{quote}
\emph{Commentary.} This is the safe-RL analogue of the energy-based
reweighting that appears in advantage-weighted regression (AWR) and KL-control
LLM finetuning. The same trick keeps appearing because it is just the
Boltzmann form of an entropy-regularised reward-max problem.

% -------------------------------------------------------------------
\section*{Paper 4 --- ``P3O'' (Penalised PPO)}
\textbf{Authors:} Linrui Zhang, Quan Vuong, Keith Ross. \textbf{Year:} 2022.
\textbf{Venue:} IJCAI. \textbf{arXiv:} 2205.11814.

\paragraph{Summary.} Replaces the saddle-point with an exact-penalty
formulation: $L_r - \kappa [\hat J_c - d]_+$. Under mild constraint
qualification, sufficiently large $\kappa$ yields a feasible minimiser of the
penalised problem that is also the constrained optimum.

\paragraph{Why it matters.} Often a one-line drop-in into PPO. The cost is
that $\kappa$ becomes a hyperparameter to tune; the benefit is that there is
no dual variable to maintain, and the loss is smoother to optimise.

\paragraph{Annotated key passage.}
\begin{quote}
``Proposition 4.1 (exact penalty). Let $\kappa^\star = \lambda^\star + \rho$
for any $\rho > 0$, where $\lambda^\star$ is the optimal Lagrange multiplier
of the original CMDP. Then for any $\kappa \geq \kappa^\star$, every local
minimiser of $L_r - \kappa[\hat J_c - d]_+$ is feasible.''
\end{quote}
\emph{Commentary.} This is a clean application of the exact-penalty theorem
of nonlinear programming. The catch: you do not know $\lambda^\star$ a
priori, so $\kappa$ is over-set in practice --- which is fine for the
penalty argument but sometimes destabilises the policy gradient.

% -------------------------------------------------------------------
\section*{Paper 5 --- ``PID-Lagrangian''}
\textbf{Authors:} Adam Stooke, Joshua Achiam, Pieter Abbeel. \textbf{Year:} 2020.
\textbf{Venue:} ICML. \textbf{arXiv:} 2007.03964.

\paragraph{Summary.} Treats the dual variable as the output of a PID
controller driven by the constraint error $\hat J_c - d$. The integral term
removes steady-state bias; the derivative term damps oscillations. Empirically
matches CPO with a much simpler implementation.

\paragraph{Why it matters.} Best-in-class \emph{practical} Lagrangian
implementation as of 2024--2025 (see the May 2025 empirical-study paper
arXiv:2510.17564, which reports PID-Lag-with-cost-normalisation as top-2
across modern Safety-Gymnasium tasks).

\paragraph{Annotated key passage.}
\begin{quote}
``We discovered that the Lagrangian's instability is best understood as a
classical control problem. The optimisation alternates between two coupled
loops (policy, multiplier) with different timescales --- the same structure
that motivates PI control in process engineering.''
\end{quote}
\emph{Commentary.} This re-framing is genuinely useful: once you start seeing
Lagrangian RL as a control problem, the bag of damping tricks
(moving averages, leaky integrators, anti-windup) becomes available
off-the-shelf. The exercise on PID-Lagrangian (workbook Ex.\ 8) sets up the
mechanics.

% -------------------------------------------------------------------
\paragraph{Sourced from.} \texttt{research/safety\_rl\_sota.md}, sections 1
(Constrained MDPs and Lagrangian methods) and 3 (risk-sensitive RL).

\end{document}
